% from main.tex
%%%%%%%%%%  scan idx 483  |  folio 470  %%%%%%%%%%
Comme $i'$ est régulière, et $p$ et $p'$ lisses, on en déduit d'après 1.3 que $i''$ et $i$ sont régulières.

La proposition qui suit ne sera pas utilisée dans la suite du Séminaire :

\textbf{Proposition 1.4.} Soit $f:X\to Y$ un morphisme plat localement de présentation finie. Pour que $f$ soit un morphisme d'intersection complète au sens de 1.1, il faut et il suffit qu'il soit un morphisme d'intersection complète au sens de EGA IV 19.3.6.

Supposons que $f$ soit un morphisme d'intersection complète au sens de 1.1, et soit $x\in X$. Alors il existe un voisinage ouvert $U$ de $x$, et une factorisation :
\[
\begin{tikzcd}[column sep=large,row sep=large]
U \arrow[r,"i"] \arrow[dr,"f"'] & V \arrow[d,"g"]\\
& Y,
\end{tikzcd}
\]
où $i$ est une immersion régulière, et $g$ un morphisme lisse. Soit $y=f(x)$ ; alors l'immersion $i_y:U_y\to V_y$ est régulière en $x$. En effet, on peut supposer, quitte à restreindre $V$, qu'il existe un $\mathcal O_V$-Module libre de type fini $E$, et un homomorphisme régulier $u:E\to\mathcal O_V$ tel que le complexe de Koszul associé à $u$ soit une résolution de $\mathcal O_U$. On considère alors l'homomorphisme $u_y$ induit sur la fibre en $y$, et le complexe de Koszul associé $K_\bullet(u_y)=K_\bullet(u)_y$. Comme $g$ est lisse, donc plat, $K_\bullet(u)$ est une résolution $Y$-plate de $\mathcal O_U$. Par suite les groupes d'homologie de $K_\bullet(u)_y$ sont les $\Tor_i^{\mathcal O_Y}(\mathcal O_U,k(y))$, où $k(y)$ est le corps résiduel de $\mathcal O_{Y,y}$. Comme $f$ est un morphisme plat, ils sont nuls, et $i_y$ est régulière en $x$. Par ailleurs, comme $g$ est lisse, l'anneau local de $V_y$ en $x$ est régulier. Comme l'anneau local de $U_y$ en $y$ est quotient d'un anneau régulier par un idéal régulier, c'est un
